\section{Setup and Notation}

Let $\mathcal{X} \subseteq \mathbb{R}^{d_0}$ denote the input space and $\mathcal{Y}$ the output space.
A BNN with $L$ layers is parameterized by:
\begin{align}
\mathbf{W} &= \{W_1, \ldots, W_L\} \quad W_l \in \{-1, +1\}^{d_l \times d_{l-1}} \\
\mathbf{b} &= \{b_1, \ldots, b_L\} \quad b_l \in \mathbb{R}^{d_l}
\end{align}

Forward pass:
\begin{align}
z^{(l)} &= W_l x^{(l-1)} + b_l \\
x^{(l)} &= \text{sign}(z^{(l)}) \in \{-1, +1\}^{d_l}
\end{align}

---

\section{Path-Sum Representation}

\subsection{Routing Pattern}
For input $x \in \mathcal{X}$, define the \emph{routing pattern} as the sequence of active (sign-determined) behaviors:
$$A(x) = \{a_1(x), \ldots, a_L(x)\} \quad a_l(x) \in \{-1, +1\}^{d_l}$$

This represents a \emph{discrete path} through the network's $2^{\sum_l d_l}$ possible neuron combinations.

\subsection{Path-Sum Model}
Express the BNN output as a weighted sum over paths:
$$f(x) = \sum_{p \in \mathcal{P}(x)} w_p \cdot g_p(x)$$
where:
\begin{itemize}
  \item $\mathcal{P}(x)$ is the set of active paths induced by routing pattern $A(x)$
  \item $w_p$ are learned weights associated with path $p$
  \item $g_p(x)$ is a path-specific function
\end{itemize}

---

\section{Key Theorem: Exponential Routing Capacity}

\subsection{Theorem 1}
Consider an $L$-layer BNN with widths $d_1, \ldots, d_L$ and binary activations. 
The number of distinct routing patterns realizable by the network is upper bounded by:
$$|R| \leq \prod_{l=1}^{L} 2^{d_l}$$

and grows \emph{exponentially in depth and width}.

\subsection{Proof Sketch}
Each binary neuron at layer $l$ contributes a binary decision: either $+1$ or $-1$.
Across all neurons in layer $l$, there are $2^{d_l}$ possible configurations.
Composing across $L$ layers, the total number of distinct routing states is at most $\prod_{l=1}^L 2^{d_l}$.

\subsection{Implication}
Even with 1-bit weights, BNNs retain \textbf{exponential expressive routing capacity}.
This explains their effectiveness despite aggressive quantization.

---

\section{Connection to SCRFP and MoE}

\subsection{Sparse Gating}
Binary activation patterns act as \emph{sparse binary gates} on features.
Each routing pattern $A(x)$ selects a specific subset of "experts" (neuron combinations).

\subsection{Equivalence}
A trained BNN with $L$ layers and width $d$ can be reinterpreted as a MoE with:
\begin{itemize}
  \item Number of experts: $\leq \prod_{l=1}^L 2^{d_l}$
  \item Gate function: Binary routing via sign activations
  \item Sparsity: Only active neurons contribute
\end{itemize}

---

\section{Visualization Strategy}

\subsection{Method 1: Layer Activation Heatmap}
For a batch of inputs, record the binary activation patterns at each layer:
$$\text{Heatmap}_{l} = [a_l(x_1) \mid a_l(x_2) \mid \cdots \mid a_l(x_N)]^T$$

This reveals:
\begin{itemize}
  \item Structured activation patterns
  \item Clustering by class
  \item Sparsity and redundancy
\end{itemize}

\subsection{Method 2: Path Overlap Graph}
Define Jaccard similarity between routing patterns:
$$\text{sim}(x_i, x_j) = \frac{|A(x_i) \cap A(x_j)|}{|A(x_i) \cup A(x_j)|}$$

Visualize as:
\begin{itemize}
  \item Adjacency matrix (heatmap)
  \item t-SNE embedding of routing vectors
  \item Network graph with inputs as nodes, similarity as edges
\end{itemize}

\subsection{Method 3: Expert Activation Histogram}
Group early-layer neurons into "expert" clusters.
Histogram activation frequency per class and per expert.

---

\section{Minimal Experiment Plan}

\subsection{Setup}
\begin{itemize}
  \item Dataset: MNIST (28×28, 10 classes)
  \item Model: 3-layer BNN (784 → 256 → 256 → 10)
  \item Training: Standard supervised learning with binarized weights
  \item Epochs: 50–100, learning rate: 0.01
\end{itemize}

\subsection{Measurements}
\begin{enumerate}
  \item Accuracy on test set
  \item Sparsity of activation patterns (fraction of active neurons)
  \item Routing similarity within classes vs. across classes
  \item Effective number of distinct routing patterns discovered
\end{enumerate}

\subsection{Outputs}
\begin{enumerate}
  \item Activation heatmaps per layer
  \item Similarity matrices / t-SNE plots
  \item Histogram of expert activations
  \item Accuracy vs. routing diversity plot
\end{enumerate}

---

\section{Expected Results}

\begin{itemize}
  \item \textbf{Strong clustering} in routing space: Samples from the same class should follow similar paths.
  \item \textbf{Sparse activation}: Few neurons active per sample (typically 30–50\% of bandwidth).
  \item \textbf{Emergent experts}: Groups of neurons specializing in class-specific routing.
  \item \textbf{Interpretability}: Visualization should clearly show discrete decision boundaries.
\end{itemize}

---

\section{Conclusion}

Binary Neural Networks are fundamentally discrete routing systems with exponential capacity in their path-sum decomposition.
By connecting BNNs to SCRFP and Mixture of Experts, we unify several quantization and modularity perspectives.
Our proposed visualizations make this structure observable and testable, enabling new interpretability and compression techniques.

\end{document}
